\begin{abstract}
Large language models memorize certain multi-token sequences as holistic units rather than composing them from parts.
We investigate whether the \futurelens method---predicting future tokens from intermediate hidden states---can detect such \noncomp phrases.
We evaluate three complementary probing signals on \gptxl: next-token probability within phrases, \logitlens decoding at nine layers spanning the full model depth, and \futurelens prediction from the first token of each phrase.
On the \farahmand noun compound dataset (1,042 compounds with expert compositionality annotations), \noncomp compounds show significantly higher within-phrase predictability than \comp ones ($p < 0.0001$, Cohen's $d = 0.49$ for surprisal), and the correlation between prediction metrics and human compositionality scores reaches Spearman $r = 0.23$ ($p < 0.0001$).
Binary classification by next-token probability achieves ROC AUC = 0.643.
On the \idiomem idiom dataset, the effect is far larger: idioms are dramatically more predictable than \comp bigrams (Cohen's $d = 1.85$).
Layer-wise analysis reveals that the \noncomp advantage emerges at layer 6 and grows through the final layers, consistent with holistic phrase processing in deeper transformer blocks.
We conclude that hidden-state prediction metrics provide a statistically significant but moderate signal for non-compositionality, best suited as a component in multi-signal detection systems rather than a standalone classifier.
\end{abstract}
